\section{The Graph This Book Builds}
\label{sec:ch05-the-graph-this-book-builds}

\index{conference example!split into three|(}%
Three services, and the split is the obvious one: whoever stores the rows owns
the type. Sessions keeps what chapter~\ref{ch:schema-design} left it. Speakers
takes the other table. Ratings is new, and it exists because a service that
contributes to a type it does not own is the arrangement this model is for, and
a two-service example never shows it.

\index{Sessions subgraph!what it declares}%
Here is the Sessions subgraph, elided as in
chapter~\ref{ch:schema-design} down to what the seam touches. Gone from it:
the \code{Query} and \code{Mutation} types and their fields, the paging types,
the \code{Node} interface, and the descriptions and \code{@cost} directives the
export writes on everything.

\begin{graphqlsdl}
schema
  @link(
    url: "https://specs.apollo.dev/federation/v2.6"
    import: ["@key", "@tag", "FieldSet"]
  ) {
  query: Query
}

type Session @key(fields: "id") {
  id: ID!
  title: String!
  abstract: String
  startsAt: DateTime!
  durationMinutes: Int!
  speaker: Speaker
}

type Speaker @key(fields: "id") {
  id: ID!
}
\end{graphqlsdl}

\index{link@\code{"@link}!the library writes it}%
The header is the library's work rather than yours, and you read it once and
then stop seeing it. It names the version of the vocabulary this document is written in, and
it imports what the document uses. The last type in the document is the whole
cost of the seam on this side: one field, so that the key has somewhere to
live.

\index{Speakers subgraph!what it declares}%
The Speakers subgraph opens with the same header and declares the same type,
and means it:

\begin{graphqlsdl}
type Speaker @key(fields: "id") {
  id: ID!
  name: String!
  bio: String
}
\end{graphqlsdl}

\index{Ratings subgraph!contributes to a type it does not own}%
And Ratings declares a type it owns and a type it does not:

\begin{graphqlsdl}
type Rating @key(fields: "id") {
  id: ID!
  score: Int!
  comment: String
  session: Session!
}

type Session @key(fields: "id") {
  id: ID!
  ratings: [Rating!]!
  averageScore: Float
}
\end{graphqlsdl}

\index{Session@\code{Session}!declared twice}%
\code{Session} now appears in two of the three documents and \code{Speaker} in
two, and in each pair one declaration is the type and the other is the key. Read
down the three and you can see the whole graph: two seams, both of them at a
\code{@key(fields: "id")}, and Sessions sitting on both.

\subsection{What the Client Gets}

\index{supergraph!hides which service answered}%
Compose the three and \code{Session} arrives as one type:

\begin{graphqlsdl}
type Session {
  id: ID!
  title: String!
  abstract: String
  startsAt: DateTime!
  durationMinutes: Int!
  speaker: Speaker
  ratings: [Rating!]!
  averageScore: Float
}
\end{graphqlsdl}

\index{supergraph!carries no federation directives}%
No \code{@key}. No second declaration. Nothing saying that two of those fields
are answered by a service which has never heard of a conference schedule. The
directives did their work at composition and are gone from what the client
reads, which is the correct outcome and the reason
chapter~\ref{ch:blame-routing} exists: a client that cannot see the seam cannot
tell you which side of it went wrong.

\index{seam!is a design decision, not a discovery}%
Look at where the two seams are, because you chose them and you could have
chosen otherwise. Sessions could have kept a speaker's name and saved a hop.
Ratings could have been a field on the Sessions service and saved a service. In
both cases the reason not to is the one chapter~\ref{ch:do-you-need-this}
started with: a seam is where one team stops needing another team's permission,
and a seam drawn anywhere else is a network call bought for nothing.

\index{seam!the bill for this one}%
The bill for this particular decomposition arrives in
chapter~\ref{ch:cross-seam-paging}. Sorting the schedule by average score is
an obvious thing for a client to ask and there is no service that can answer
it: Sessions owns the ordering and knows no scores, Ratings owns the scores
and knows no schedule, and no directive in
section~\ref{sec:ch05-the-directives} closes that gap. I chose this example
partly because that problem is unavoidable in it rather than contrived.

\subsection{What Part II Does With It}

\index{Part II!the four chapters that build this}%
Chapter~\ref{ch:first-subgraph} makes the service you already have into a
subgraph: the package, the key, the reference resolver, and the two fields the
contract adds to \code{Query}. It gets most of the way to the first of those
three documents and stops short of one thing. Sessions still stores the
speakers, so its \code{Speaker} is the whole type rather than the stub above,
and it stays that way until there is a service to move the rows to. That
chapter ships one subgraph and composes nothing.

Chapter~\ref{ch:composition} composes, which means meeting the composer's
opinions about the schema you wrote, including the one
section~\ref{sec:ch05-who-owns-a-field} promised you about \code{Query.node}.
Chapter~\ref{ch:enter-the-router} puts a router in front and answers the first
query that crosses a seam. Chapter~\ref{ch:second-third-subgraph} adds the other
two services and the directives in
section~\ref{sec:ch05-the-directives}'s seam group, and ends with the graph
part~III spends nine chapters breaking.
\index{conference example!split into three|)}%
